\documentclass[../main.tex]{subfiles}
\usepackage{amsmath}
\usepackage{hyperref}
\usepackage{enumitem}

\begin{document}

This project has provided a comprehensive exploration of adversarial attacks and defenses in deep learning, spanning both computer vision and language processing domains. By systematically examining a wide spectrum of tasks—including image classification, segmentation, object detection, natural language processing, large language models, and automatic speech recognition—we have highlighted the pervasive vulnerability of modern neural networks to carefully crafted, often imperceptible, adversarial perturbations. Our study underscores that these vulnerabilities are not isolated to specific architectures or data modalities, but rather represent a fundamental challenge to the reliability and security of AI systems as they are increasingly deployed in safety-critical and real-world applications.

Throughout this work, we have emphasized the importance of evaluating both attack strategies and defense mechanisms in a unified, cross-domain framework. This approach enables a deeper understanding of the trade-offs between robustness, accuracy, and computational efficiency, and reveals that effective protection against adversarial threats typically requires adaptive, task-specific solutions rather than one-size-fits-all defenses. By bridging empirical analysis with practical mitigation strategies, this project lays the groundwork for more resilient and trustworthy AI, informing future research and deployment in domains where security and reliability are paramount, such as autonomous vehicles, healthcare, and secure communications.

\end{document}